\documentclass[11pt,a4paper]{article}
\usepackage[utf8]{inputenc}
\usepackage[T1]{fontenc}
\usepackage[spanish,es-minimal,es-noshorthands]{babel}
\AtBeginDocument{\renewcommand{\sin}{\operatorname{sen}}\renewcommand{\arcsin}{\operatorname{arcsen}}}
\usepackage{lmodern}
\usepackage[sfdefault,lining]{FiraSans}
\usepackage[a4paper,top=2.3cm,bottom=2.2cm,left=1.8cm,right=1.8cm,headheight=1.3cm,headsep=0.4cm,footskip=1cm]{geometry}
\usepackage{amsmath,amssymb}
\usepackage{icomma}
\usepackage{graphicx}
\usepackage[table]{xcolor}
\usepackage{tikz}
\usepackage{fontawesome5}
\usepackage{ccicons}
\usepackage{fancyhdr}
\usepackage{lastpage}
\usepackage{tasks}
\usepackage{needspace}
\usepackage[most]{tcolorbox}
\usepackage{array}
\usepackage[hidelinks]{hyperref}
\AtBeginDocument{\hypersetup{pdftitle={Soluciones de la ficha de la unidad 10: Distribuciones de probabilidad. 1º Bach.},pdfsubject={Matemáticas CCSS 1º Bach.},pdfkeywords={pdfkeywords={variable aleatoria, función de probabilidad, función de distribución, esperanza, distribución binomial, función de densidad, distribución uniforme, distribución normal, normal tipificada, tipificación, aproximación de la binomial a la normal}}}}

\definecolor{navy}{HTML}{1B2748}
\definecolor{accent}{HTML}{E8603A}
\definecolor{teal}{HTML}{1F9E90}
\definecolor{soft}{HTML}{EEF1F8}
\definecolor{muted}{HTML}{5D6685}

\setlength{\parindent}{0pt}
\setlength{\parskip}{3pt}
\settasks{label=\alph*), label-format=\bfseries\color{navy}, label-width=1.3em, item-indent=1.8em, column-sep=0.6em, after-item-skip=5pt, before-skip=3pt, after-skip=4pt}

% ---------- capçalera i peu ----------
\newcommand{\llicencia}{\href{https://creativecommons.org/licenses/by-nc-sa/4.0/deed.es}{\ccbyncsa}\ \ Licencia Creative Commons Atribución-NoComercial-CompartirIgual 4.0 Internacional (CC BY-NC-SA 4.0)}
\pagestyle{fancy}
\fancyhf{}
\fancyhead[L]{\raisebox{-0.35\height}{\includegraphics[height=0.95cm]{logo-gva.png}}}
\fancyhead[C]{\small\color{muted}Soluciones de la ficha · Unidad 10: Distribuciones de probabilidad · 1\textsuperscript{o} Bach.}
\fancyhead[R]{\raisebox{-0.35\height}{\includegraphics[height=0.95cm]{logo-vermellar.png}}}
\fancyfoot[L]{\scriptsize\color{muted}\llicencia}
\fancyfoot[R]{\small\color{muted}\thepage}
\renewcommand{\headrulewidth}{0.4pt}
\renewcommand{\headrule}{\color{muted!40}\hrule height \headrulewidth}
\fancypagestyle{primera}{\fancyhf{}\renewcommand{\headrulewidth}{0pt}
  \fancyfoot[L]{\scriptsize\color{muted}\llicencia}
  \fancyfoot[R]{\small\color{muted}\thepage}}

% ---------- nivells i exercicis ----------
\newcommand{\estrelles}[1]{{\color{accent}\ifcase#1\or\faStar\faStar[regular]\faStar[regular]\or\faStar\faStar\faStar[regular]\or\faStar\faStar\faStar\fi}}
\newcounter{exer}
\newcommand{\ex}[1]{\par\needspace{4\baselineskip}\medskip\refstepcounter{exer}\noindent{\color{navy}\bfseries\theexer.}\ \ }

\newcommand{\seccio}[2]{\par\needspace{16\baselineskip}\bigskip
  \begin{tcolorbox}[colback=navy,colframe=navy,coltext=white,arc=2mm,boxrule=0pt,left=3mm,right=3mm,top=1.2mm,bottom=1.2mm]
  {\large\bfseries #1\quad #2}
  \end{tcolorbox}\vspace{-2pt}}

\usepackage{enumitem}
\newcommand{\exf}[1]{\par\medskip\refstepcounter{exer}\noindent{\color{navy}\bfseries\theexer.}\ \ }
\providecommand{\gt}{>}\providecommand{\lt}{<}
% quadre de recordatori (només a la fitxa d'exercicis)
\newcommand{\recorda}[1]{\par\smallskip\begin{tcolorbox}[colback=accent!7,colframe=accent!55,arc=2mm,boxrule=0.5pt,left=2.5mm,right=2.5mm,top=1mm,bottom=1mm]\small{\color{accent}\faLightbulb[regular]}\ \textbf{Recuerda:}\ #1\end{tcolorbox}}
\newcommand{\pau}[1]{{\color{teal}\small\bfseries\faUniversity\ PAU · #1.}\ }
\newcommand{\res}[1]{\boxed{#1}}\newcommand{\rest}[1]{\fbox{#1}}

\def\fmneed{11}% que el títol d'una secció no quede sol al final del full
% =====================================================================
% Capçalera FentMates de les fitxes: un color per a cada curs.
% Els colors són els de la pàgina inicial (LEVELS a assets/js/materials.js);
% el logo ve de _eines/logo/pdf/ (python3 _eines/logo/genera.py).
%
% Ús, al final del preàmbul de la unitat:
%   % =====================================================================
% Capçalera FentMates de les fitxes: un color per a cada curs.
% Els colors són els de la pàgina inicial (LEVELS a assets/js/materials.js);
% el logo ve de _eines/logo/pdf/ (python3 _eines/logo/genera.py).
%
% Ús, al final del preàmbul de la unitat:
%   % =====================================================================
% Capçalera FentMates de les fitxes: un color per a cada curs.
% Els colors són els de la pàgina inicial (LEVELS a assets/js/materials.js);
% el logo ve de _eines/logo/pdf/ (python3 _eines/logo/genera.py).
%
% Ús, al final del preàmbul de la unitat:
%   \input{../comu/fentmates.tex}
%   \fmcurs{1eso}
% i, al començament del document:
%   \capcalera{Matemàtiques}{1r ESO}{Fitxa d'exercicis · Unitat 2}{Divisibilitat}
% Opcional abans de \input: \def\fmneed{14} (espai mínim abans de cada secció).
% =====================================================================
\makeatletter
\def\fm@logodir{../../logo/pdf}
\providecommand{\fmneed}{10}
\def\fm@cursnom{}

% clau del curs, color i color del text damunt del color
\newcommand{\fm@curs}[3]{\@namedef{fm@color@#1}{#2}\@namedef{fm@text@#1}{#3}}
\fm@curs{1eso}{E6B84A}{navy}
\fm@curs{2eso}{3B82F6}{white}
\fm@curs{3eso}{FF7A50}{white}
\fm@curs{4eso-a}{A78BFA}{white}
\fm@curs{4eso-b}{6D28D9}{white}
\fm@curs{1batx-mi}{9B1C1C}{white}
\fm@curs{1batx-ccss}{EF6B6B}{white}
\fm@curs{2batx-mii}{1E6B47}{white}
\fm@curs{2batx-ccss}{2BB3A3}{white}

\newcommand{\fmcurs}[1]{%
  \@ifundefined{fm@color@#1}{\PackageError{fentmates}{Curs desconegut: #1}{}}{}%
  \edef\fm@tmp{\noexpand\definecolor{curs}{HTML}{\@nameuse{fm@color@#1}}%
               \noexpand\colorlet{curstext}{\@nameuse{fm@text@#1}}}\fm@tmp
  \colorlet{cursfosc}{curs!72!black}%
  \def\fm@logo{\fm@logodir/vertical-#1.pdf}%
  \def\fm@marca{\fm@logodir/marca-#1.pdf}}

% ---------- capçalera de la primera pàgina ----------
% #1 matèria (petit), #2 curs, #3 tipus de document i unitat, #4 títol
\newcommand{\capcalera}[4]{%
  \gdef\fm@cursnom{#2}%
  \par\noindent
  \begin{tcolorbox}[colback=curs,colframe=curs,boxrule=0pt,arc=2.2mm,width=3.5cm,height=2.15cm,
      valign=center,left=3mm,right=2mm,top=1mm,bottom=1mm,nobeforeafter,box align=center]
    \color{curstext}{\footnotesize #1}\par\vspace{1pt}{\fontsize{19}{21}\selectfont\bfseries #2\par}%
  \end{tcolorbox}\hspace{5mm}%
  \begin{minipage}[c]{\dimexpr\linewidth-3.5cm-5mm-2.9cm\relax}
    {\color{cursfosc}\small\bfseries\MakeUppercase{#3}}\par\vspace{2pt}
    {\color{navy}\fontsize{20}{23}\selectfont\bfseries #4\par}
  \end{minipage}\hfill
  \begin{minipage}[c]{2.6cm}\raggedleft\includegraphics[width=2.6cm]{\fm@logo}\end{minipage}\par
  \vspace{2pt}}

% ---------- barres de secció del color del curs ----------
\renewcommand{\seccio}[2]{\par\needspace{\fmneed\baselineskip}\bigskip
  \begin{tcolorbox}[colback=curs,colframe=curs,coltext=curstext,arc=2mm,boxrule=0pt,left=3mm,right=3mm,top=1.2mm,bottom=1.2mm]
  {\large\bfseries #1\quad #2}
  \end{tcolorbox}\vspace{-2pt}}

% ---------- estrelles de dificultat del color del curs (a Batxillerat no se'n posen) ----------
\@ifundefined{estrelles}{}{%
  \renewcommand{\estrelles}[1]{{\color{curs}\ifcase#1\or\faStar\faStar[regular]\faStar[regular]\or\faStar\faStar\faStar[regular]\or\faStar\faStar\faStar\fi}}}

% ---------- capçalera i peu de les pàgines ----------
% primera pàgina: logos institucionals al peu
\fancypagestyle{primera}{\fancyhf{}\renewcommand{\headrulewidth}{0pt}
  \fancyfoot[L]{\raisebox{-0.32\height}{\includegraphics[height=0.78cm]{logo-gva.png}}\hspace{5mm}%
    \raisebox{-0.32\height}{\includegraphics[height=0.78cm]{logo-vermellar.png}}}
  \fancyfoot[R]{\scriptsize\color{muted}\href{https://creativecommons.org/licenses/by-nc-sa/4.0/deed.ca}{\ccbyncsa}\ \ CC BY-NC-SA 4.0 · fentmates.org\qquad\small\thepage}}
% la resta: el símbol a l'esquerra i el curs a la dreta, amb el títol al centre com abans
\fancyhead[L]{\raisebox{-0.3\height}{\includegraphics[height=0.62cm]{\fm@marca}}}
\fancyhead[R]{\small\bfseries\color{cursfosc}\fm@cursnom}
\renewcommand{\headrule}{\color{curs!45}\hrule height \headrulewidth}
\makeatother

%   \fmcurs{1eso}
% i, al començament del document:
%   \capcalera{Matemàtiques}{1r ESO}{Fitxa d'exercicis · Unitat 2}{Divisibilitat}
% Opcional abans de \input: \def\fmneed{14} (espai mínim abans de cada secció).
% =====================================================================
\makeatletter
\def\fm@logodir{../../logo/pdf}
\providecommand{\fmneed}{10}
\def\fm@cursnom{}

% clau del curs, color i color del text damunt del color
\newcommand{\fm@curs}[3]{\@namedef{fm@color@#1}{#2}\@namedef{fm@text@#1}{#3}}
\fm@curs{1eso}{E6B84A}{navy}
\fm@curs{2eso}{3B82F6}{white}
\fm@curs{3eso}{FF7A50}{white}
\fm@curs{4eso-a}{A78BFA}{white}
\fm@curs{4eso-b}{6D28D9}{white}
\fm@curs{1batx-mi}{9B1C1C}{white}
\fm@curs{1batx-ccss}{EF6B6B}{white}
\fm@curs{2batx-mii}{1E6B47}{white}
\fm@curs{2batx-ccss}{2BB3A3}{white}

\newcommand{\fmcurs}[1]{%
  \@ifundefined{fm@color@#1}{\PackageError{fentmates}{Curs desconegut: #1}{}}{}%
  \edef\fm@tmp{\noexpand\definecolor{curs}{HTML}{\@nameuse{fm@color@#1}}%
               \noexpand\colorlet{curstext}{\@nameuse{fm@text@#1}}}\fm@tmp
  \colorlet{cursfosc}{curs!72!black}%
  \def\fm@logo{\fm@logodir/vertical-#1.pdf}%
  \def\fm@marca{\fm@logodir/marca-#1.pdf}}

% ---------- capçalera de la primera pàgina ----------
% #1 matèria (petit), #2 curs, #3 tipus de document i unitat, #4 títol
\newcommand{\capcalera}[4]{%
  \gdef\fm@cursnom{#2}%
  \par\noindent
  \begin{tcolorbox}[colback=curs,colframe=curs,boxrule=0pt,arc=2.2mm,width=3.5cm,height=2.15cm,
      valign=center,left=3mm,right=2mm,top=1mm,bottom=1mm,nobeforeafter,box align=center]
    \color{curstext}{\footnotesize #1}\par\vspace{1pt}{\fontsize{19}{21}\selectfont\bfseries #2\par}%
  \end{tcolorbox}\hspace{5mm}%
  \begin{minipage}[c]{\dimexpr\linewidth-3.5cm-5mm-2.9cm\relax}
    {\color{cursfosc}\small\bfseries\MakeUppercase{#3}}\par\vspace{2pt}
    {\color{navy}\fontsize{20}{23}\selectfont\bfseries #4\par}
  \end{minipage}\hfill
  \begin{minipage}[c]{2.6cm}\raggedleft\includegraphics[width=2.6cm]{\fm@logo}\end{minipage}\par
  \vspace{2pt}}

% ---------- barres de secció del color del curs ----------
\renewcommand{\seccio}[2]{\par\needspace{\fmneed\baselineskip}\bigskip
  \begin{tcolorbox}[colback=curs,colframe=curs,coltext=curstext,arc=2mm,boxrule=0pt,left=3mm,right=3mm,top=1.2mm,bottom=1.2mm]
  {\large\bfseries #1\quad #2}
  \end{tcolorbox}\vspace{-2pt}}

% ---------- estrelles de dificultat del color del curs (a Batxillerat no se'n posen) ----------
\@ifundefined{estrelles}{}{%
  \renewcommand{\estrelles}[1]{{\color{curs}\ifcase#1\or\faStar\faStar[regular]\faStar[regular]\or\faStar\faStar\faStar[regular]\or\faStar\faStar\faStar\fi}}}

% ---------- capçalera i peu de les pàgines ----------
% primera pàgina: logos institucionals al peu
\fancypagestyle{primera}{\fancyhf{}\renewcommand{\headrulewidth}{0pt}
  \fancyfoot[L]{\raisebox{-0.32\height}{\includegraphics[height=0.78cm]{logo-gva.png}}\hspace{5mm}%
    \raisebox{-0.32\height}{\includegraphics[height=0.78cm]{logo-vermellar.png}}}
  \fancyfoot[R]{\scriptsize\color{muted}\href{https://creativecommons.org/licenses/by-nc-sa/4.0/deed.ca}{\ccbyncsa}\ \ CC BY-NC-SA 4.0 · fentmates.org\qquad\small\thepage}}
% la resta: el símbol a l'esquerra i el curs a la dreta, amb el títol al centre com abans
\fancyhead[L]{\raisebox{-0.3\height}{\includegraphics[height=0.62cm]{\fm@marca}}}
\fancyhead[R]{\small\bfseries\color{cursfosc}\fm@cursnom}
\renewcommand{\headrule}{\color{curs!45}\hrule height \headrulewidth}
\makeatother

%   \fmcurs{1eso}
% i, al començament del document:
%   \capcalera{Matemàtiques}{1r ESO}{Fitxa d'exercicis · Unitat 2}{Divisibilitat}
% Opcional abans de \input: \def\fmneed{14} (espai mínim abans de cada secció).
% =====================================================================
\makeatletter
\def\fm@logodir{../../logo/pdf}
\providecommand{\fmneed}{10}
\def\fm@cursnom{}

% clau del curs, color i color del text damunt del color
\newcommand{\fm@curs}[3]{\@namedef{fm@color@#1}{#2}\@namedef{fm@text@#1}{#3}}
\fm@curs{1eso}{E6B84A}{navy}
\fm@curs{2eso}{3B82F6}{white}
\fm@curs{3eso}{FF7A50}{white}
\fm@curs{4eso-a}{A78BFA}{white}
\fm@curs{4eso-b}{6D28D9}{white}
\fm@curs{1batx-mi}{9B1C1C}{white}
\fm@curs{1batx-ccss}{EF6B6B}{white}
\fm@curs{2batx-mii}{1E6B47}{white}
\fm@curs{2batx-ccss}{2BB3A3}{white}

\newcommand{\fmcurs}[1]{%
  \@ifundefined{fm@color@#1}{\PackageError{fentmates}{Curs desconegut: #1}{}}{}%
  \edef\fm@tmp{\noexpand\definecolor{curs}{HTML}{\@nameuse{fm@color@#1}}%
               \noexpand\colorlet{curstext}{\@nameuse{fm@text@#1}}}\fm@tmp
  \colorlet{cursfosc}{curs!72!black}%
  \def\fm@logo{\fm@logodir/vertical-#1.pdf}%
  \def\fm@marca{\fm@logodir/marca-#1.pdf}}

% ---------- capçalera de la primera pàgina ----------
% #1 matèria (petit), #2 curs, #3 tipus de document i unitat, #4 títol
\newcommand{\capcalera}[4]{%
  \gdef\fm@cursnom{#2}%
  \par\noindent
  \begin{tcolorbox}[colback=curs,colframe=curs,boxrule=0pt,arc=2.2mm,width=3.5cm,height=2.15cm,
      valign=center,left=3mm,right=2mm,top=1mm,bottom=1mm,nobeforeafter,box align=center]
    \color{curstext}{\footnotesize #1}\par\vspace{1pt}{\fontsize{19}{21}\selectfont\bfseries #2\par}%
  \end{tcolorbox}\hspace{5mm}%
  \begin{minipage}[c]{\dimexpr\linewidth-3.5cm-5mm-2.9cm\relax}
    {\color{cursfosc}\small\bfseries\MakeUppercase{#3}}\par\vspace{2pt}
    {\color{navy}\fontsize{20}{23}\selectfont\bfseries #4\par}
  \end{minipage}\hfill
  \begin{minipage}[c]{2.6cm}\raggedleft\includegraphics[width=2.6cm]{\fm@logo}\end{minipage}\par
  \vspace{2pt}}

% ---------- barres de secció del color del curs ----------
\renewcommand{\seccio}[2]{\par\needspace{\fmneed\baselineskip}\bigskip
  \begin{tcolorbox}[colback=curs,colframe=curs,coltext=curstext,arc=2mm,boxrule=0pt,left=3mm,right=3mm,top=1.2mm,bottom=1.2mm]
  {\large\bfseries #1\quad #2}
  \end{tcolorbox}\vspace{-2pt}}

% ---------- estrelles de dificultat del color del curs (a Batxillerat no se'n posen) ----------
\@ifundefined{estrelles}{}{%
  \renewcommand{\estrelles}[1]{{\color{curs}\ifcase#1\or\faStar\faStar[regular]\faStar[regular]\or\faStar\faStar\faStar[regular]\or\faStar\faStar\faStar\fi}}}

% ---------- capçalera i peu de les pàgines ----------
% primera pàgina: logos institucionals al peu
\fancypagestyle{primera}{\fancyhf{}\renewcommand{\headrulewidth}{0pt}
  \fancyfoot[L]{\raisebox{-0.32\height}{\includegraphics[height=0.78cm]{logo-gva.png}}\hspace{5mm}%
    \raisebox{-0.32\height}{\includegraphics[height=0.78cm]{logo-vermellar.png}}}
  \fancyfoot[R]{\scriptsize\color{muted}\href{https://creativecommons.org/licenses/by-nc-sa/4.0/deed.ca}{\ccbyncsa}\ \ CC BY-NC-SA 4.0 · fentmates.org\qquad\small\thepage}}
% la resta: el símbol a l'esquerra i el curs a la dreta, amb el títol al centre com abans
\fancyhead[L]{\raisebox{-0.3\height}{\includegraphics[height=0.62cm]{\fm@marca}}}
\fancyhead[R]{\small\bfseries\color{cursfosc}\fm@cursnom}
\renewcommand{\headrule}{\color{curs!45}\hrule height \headrulewidth}
\makeatother

\fmcurs{1batx-ccss}

\begin{document}
\thispagestyle{primera}

\capcalera{Matemáticas CCSS}{1\textsuperscript{o} Bach.}{Soluciones de la ficha · Unidad 10}{\raggedright Distribuciones de probabilidad}
\vspace{8pt}

\begin{tcolorbox}[colback=soft,colframe=soft,arc=2mm,left=3mm,right=3mm,top=1.5mm,bottom=1.5mm]
\small Las probabilidades de la normal se han calculado redondeando $z$ a dos decimales y leyendo la tabla $N(0,1)$, como se hace en la PAU. En los problemas inversos se ha tomado el valor de la tabla más próximo; por eso algunos resultados pueden diferir ligeramente de los que se obtienen con una calculadora.
\end{tcolorbox}

% =====================================================
\seccio{1}{Variables aleatorias discretas}

\exf{1} \begin{tasks}(2)
  \task \rest{Discreta}
  \task \rest{Continua}
  \task \rest{Continua}
  \task \rest{Discreta}
  \task \rest{Continua}
  \task \rest{Discreta}
\end{tasks}

\ex{1} \begin{tasks}(1)
  \task Espacio muestral: CC, CX, XC, XX. \ $P(X=0)=\frac14$, $P(X=1)=\frac12$, $P(X=2)=\frac14$.

\begin{center}\begin{tikzpicture}[x=0.8cm,y=4.00cm,baseline=(current bounding box.center)]
\draw[muted,-latex] (-0.70,0)--(2.80,0);
\fill[accent!70] (-0.12,0) rectangle (0.12,0.2500); \node[below,font=\scriptsize,text=navy] at (0.00,0) {$0$}; \node[above,font=\scriptsize,text=navy] at (0.00,0.2500) {$\frac{1}{4}$};
\fill[accent!70] (0.88,0) rectangle (1.12,0.5000); \node[below,font=\scriptsize,text=navy] at (1.00,0) {$1$}; \node[above,font=\scriptsize,text=navy] at (1.00,0.5000) {$\frac{1}{2}$};
\fill[accent!70] (1.88,0) rectangle (2.12,0.2500); \node[below,font=\scriptsize,text=navy] at (2.00,0) {$2$}; \node[above,font=\scriptsize,text=navy] at (2.00,0.2500) {$\frac{1}{4}$};
\end{tikzpicture}\end{center}

  \task $F(0)=\frac14$, \ $F(1)=\frac34$, \ $F(2)=1$.
  \task $\mu=0+\frac12+\frac24=\res{1}$; \ $\sigma^2=\frac12+1-1=\frac12\Rightarrow\res{\sigma=\sqrt{0{,}5}\approx0{,}707}$
\end{tasks}

\ex{1} \begin{tasks}(1)
  \task $k+2k+3k+4k=1\Rightarrow10k=1\Rightarrow\res{k=0{,}1}$
  \task $\mu=0{,}1+0{,}4+0{,}9+1{,}6=\res{3}$; \ $\sigma^2=0{,}1+0{,}8+2{,}7+6{,}4-3^2=1\Rightarrow\res{\sigma=1}$
  \task $P(X\geq3)=0{,}3+0{,}4=\res{0{,}7}$; \ $P(X<3)=1-0{,}7=\res{0{,}3}$
\end{tasks}

\ex{2} \par{\centering\small\renewcommand{\arraystretch}{1.3}\begin{tabular}{c|ccccccccccc}$x_i$ & 2 & 3 & 4 & 5 & 6 & 7 & 8 & 9 & 10 & 11 & 12\\\hline $p_i$ & $\frac{1}{36}$ & $\frac{2}{36}$ & $\frac{3}{36}$ & $\frac{4}{36}$ & $\frac{5}{36}$ & $\frac{6}{36}$ & $\frac{5}{36}$ & $\frac{4}{36}$ & $\frac{3}{36}$ & $\frac{2}{36}$ & $\frac{1}{36}$\end{tabular}\par}\smallskip
$P(X\geq10)=\frac{3+2+1}{36}=\res{\tfrac16}$; \ $\mu=\frac{252}{36}=\res{7}$ (por simetría); \ $\sigma^2=\frac{1974}{36}-49=\frac{35}{6}\Rightarrow\res{\sigma\approx2{,}415}$.

\ex{2} $P(X=0)=\frac46\cdot\frac35=\frac25$, \ $P(X=1)=\frac46\cdot\frac25+\frac26\cdot\frac45=\frac8{15}$, \ $P(X=2)=\frac26\cdot\frac15=\frac1{15}$. \ $\mu=\frac8{15}+\frac2{15}=\res{\tfrac23}$; \ $\sigma^2=\frac8{15}+\frac4{15}-\frac49=\frac{16}{45}\Rightarrow\res{\sigma\approx0{,}596}$.

\ex{2} $X$ toma los valores $3$, $1$ y $-1$ con probabilidades $\frac16$, $\frac16$ y $\frac46$. \ $\mu=\frac{3+1-4}{6}=\res{0}$: \rest{es un juego justo}. \ $\sigma^2=\frac{9+1+4}{6}=\frac73\Rightarrow\res{\sigma\approx1{,}53}$ euros.

\ex{3} $\left\{\begin{array}{l}a+0{,}5+b=1\\ 0{,}5+2b=1{,}2\end{array}\right.\Rightarrow\res{b=0{,}35},\ \res{a=0{,}15}$; \ $\sigma^2=0{,}5+4\cdot0{,}35-1{,}2^2=\res{0{,}46}$.

\ex{2} Restamos valores consecutivos: $p_0=0{,}1$, $p_1=0{,}3$, $p_2=0{,}4$, $p_3=0{,}2$. \ $P(1\leq X\leq2)=F(2)-F(0)=\res{0{,}7}$; \ $\mu=0{,}3+0{,}8+0{,}6=\res{1{,}7}$.


% =====================================================
\seccio{2}{La distribución binomial}

\exf{1} \begin{tasks}(1)
  \task \rest{Sí}: $B\left(12,\frac16\right)$.
  \task \rest{No}: sin reemplazamiento, la probabilidad de roja cambia en cada extracción.
  \task \rest{Sí}: $B(50;\ 0{,}02)$ (la producción es muy grande: las pruebas son prácticamente independientes).
  \task \rest{No}: el número de pruebas no está fijado.
\end{tasks}

\ex{1} $X\sim B\left(5,\frac16\right)$: \ $P(X=2)=\binom52\left(\frac16\right)^2\left(\frac56\right)^3=10\cdot\frac{125}{7776}=\res{0{,}1608}$; \ $P(X\geq1)=1-\left(\frac56\right)^5=\res{0{,}5981}$.

\ex{2} \begin{tasks}(1)
  \task $X\sim B\left(8,\frac13\right)$; \ $\mu=np=\frac83\approx\res{2{,}67}$ aciertos; \ $\sigma=\sqrt{8\cdot\frac13\cdot\frac23}=\frac43\approx\res{1{,}33}$
  \task $P(X=4)=\binom84\left(\frac13\right)^4\left(\frac23\right)^4=70\cdot\frac{16}{6561}=\res{0{,}1707}$; \ $P(X\geq1)=1-\left(\frac23\right)^8=\res{0{,}9610}$
\end{tasks}

\ex{2} $X=$ «número de clientes que no compran en línea» $\sim B(10;\ 0{,}1)$.\begin{tasks}(1)
  \task $P(X=2)=\binom{10}{2}0{,}1^2\,0{,}9^8=\res{0{,}1937}$
  \task $P(X<2)=P(X=0)+P(X=1)=0{,}9^{10}+10\cdot0{,}1\cdot0{,}9^9=\res{0{,}7361}$
\end{tasks}

\ex{1} Bloque $n=5$, columna $p=0{,}25$:\begin{tasks}(2)
  \task $\res{0{,}2637}$
  \task $0{,}2373+0{,}3955=\res{0{,}6328}$
  \task $0{,}0146+0{,}0010=\res{0{,}0156}$
  \task $1-0{,}2373=\res{0{,}7627}$
\end{tasks}

\ex{2} $X\sim B(7;\ 0{,}8)$ no está en la tabla ($p>0{,}5$). Contamos los no satisfechos: $Y=7-X\sim B(7;\ 0{,}2)$.\begin{tasks}(1)
  \task $P(X=6)=P(Y=1)=\res{0{,}3670}$
  \task $P(X\geq6)=P(Y\leq1)=0{,}2097+0{,}3670=\res{0{,}5767}$
  \task $P(X\leq5)=1-P(X\geq6)=\res{0{,}4233}$
\end{tasks}

\ex{1} $X=$ «número de chicos» $\sim B(4;\ 0{,}5)$: \ $P(X=2)=\binom42\,0{,}5^4=\res{0{,}3750}$; \ $P(\text{al menos una chica})=1-P(X=4)=1-0{,}0625=\res{0{,}9375}$; \ $\mu=4\cdot0{,}5=\res{2}$ chicos.

\ex{2} \begin{tasks}(1)
  \task $X\sim B(50;\ 0{,}03)$: \ $\mu=50\cdot0{,}03=\res{1{,}5}$; \ $\sigma=\sqrt{50\cdot0{,}03\cdot0{,}97}\approx\res{1{,}21}$
  \task $P(X=0)=0{,}97^{50}=\res{0{,}2181}$; \ $P(X\leq1)=0{,}97^{50}+50\cdot0{,}03\cdot0{,}97^{49}=\res{0{,}5553}$
\end{tasks}

\ex{3} $P(\text{al menos una cara})=1-0{,}5^n>0{,}99\iff0{,}5^n<0{,}01\iff2^n>100$. Como $2^6=64$ y $2^7=128$: \ $\res{n=7}$ lanzamientos.


% =====================================================
\seccio{3}{Variables aleatorias continuas. Función de densidad}

\exf{1} \begin{tasks}(1)
  \task $f(x)\geq0$ y el área es un rectángulo de base 6 y altura $\frac16$: $6\cdot\frac16=1$ \checkmark

\begin{center}\begin{tikzpicture}[x=0.500cm,y=6.000cm,baseline=(current bounding box.center)]
\fill[accent!35] (3.000,0) -- (3.000,0.1667) -- (7.000,0.1667) -- (7.000,0) -- cycle;
\draw[muted,-latex] (1.20,0)--(8.80,0) node[right,font=\scriptsize]{$x$};
\draw[muted,-latex] (1.50,0)--(1.50,0.217);
\draw[teal,thick] (1.20,0) -- (2.000,0) -- (2.000,0.1667) -- (8.000,0.1667) -- (8.000,0) -- (8.50,0);
\draw[muted] (2.000,0.03)--(2.000,-0.03) node[below,font=\scriptsize,text=navy]{$2$};
\draw[muted] (3.000,0.03)--(3.000,-0.03) node[below,font=\scriptsize,text=navy]{$3$};
\draw[muted] (7.000,0.03)--(7.000,-0.03) node[below,font=\scriptsize,text=navy]{$7$};
\draw[muted] (8.000,0.03)--(8.000,-0.03) node[below,font=\scriptsize,text=navy]{$8$};
\draw[muted,dashed] (1.50,0.1667)--(8.000,0.1667); \node[left,font=\scriptsize,text=navy] at (1.50,0.1667) {$\frac16$};
\end{tikzpicture}\end{center}

  \task $P(X<4)=2\cdot\frac16=\res{\tfrac13}$; \ $P(3\leq X\leq7)=4\cdot\frac16=\res{\tfrac23}$; \ $P(X=5)=\res{0}$
  \task $\mu=\frac{2+8}{2}=\res{5}$
\end{tasks}

\ex{2} La región es un triángulo de base 4 y altura $f(4)=4k$: $\frac{4\cdot4k}{2}=8k=1\Rightarrow\res{k=\tfrac18}$. \ $P(X\leq2)$: triángulo de base 2 y altura $f(2)=\frac14$: $\frac{2\cdot\frac14}{2}=\res{\tfrac14}$. \ $P(X>3)=1-P(X\leq3)=1-\frac{3\cdot\frac38}{2}=1-\frac9{16}=\res{\tfrac7{16}}$.

\begin{center}\begin{tikzpicture}[x=0.700cm,y=3.000cm,baseline=(current bounding box.center)]
\fill[accent!35] (3.000,0) -- (3.000,0.3750) -- (4.000,0.5000) -- (4.000,0) -- cycle;
\draw[muted,-latex] (-0.80,0)--(4.80,0) node[right,font=\scriptsize]{$x$};
\draw[muted,-latex] (-0.50,0)--(-0.50,0.650);
\draw[teal,thick] (-0.80,0) -- (0.000,0) -- (0.000,0.0000) -- (4.000,0.5000) -- (4.000,0) -- (4.50,0);
\draw[muted] (0.000,0.03)--(0.000,-0.03) node[below,font=\scriptsize,text=navy]{$0$};
\draw[muted] (2.000,0.03)--(2.000,-0.03) node[below,font=\scriptsize,text=navy]{$2$};
\draw[muted] (3.000,0.03)--(3.000,-0.03) node[below,font=\scriptsize,text=navy]{$3$};
\draw[muted] (4.000,0.03)--(4.000,-0.03) node[below,font=\scriptsize,text=navy]{$4$};
\draw[muted,dashed] (-0.50,0.5000)--(4.000,0.5000); \node[left,font=\scriptsize,text=navy] at (-0.50,0.5000) {$\frac12$};
\end{tikzpicture}\end{center}


\ex{1} $X\sim U(0,20)$, $f(x)=\frac1{20}$: \ $P(X>15)=\frac5{20}=\res{0{,}25}$; \ $P(5\leq X\leq12)=\frac7{20}=\res{0{,}35}$; \ $\mu=\res{10}$ minutos.

\ex{2} \begin{tasks}(1)
  \task $\frac{4\cdot h}{2}=1\Rightarrow\res{h=\tfrac12}$. La recta de la izquierda es $f(x)=\frac x4$.
  \task $P(X\leq1)$: triángulo de base 1 y altura $f(1)=\frac14$: $\res{\tfrac18}$; \ por simetría, $P(1\leq X\leq3)=1-2\cdot\frac18=\res{\tfrac34}$; \ $P(X\geq3{,}5)$: triángulo de base $0{,}5$ y altura $\frac18$: $\res{\tfrac1{32}}$

\begin{center}\begin{tikzpicture}[x=0.700cm,y=3.000cm,baseline=(current bounding box.center)]
\fill[accent!35] (1.000,0) -- (1.000,0.2500) -- (2.000,0.5000) -- (3.000,0.2500) -- (3.000,0) -- cycle;
\draw[muted,-latex] (-0.80,0)--(4.80,0) node[right,font=\scriptsize]{$x$};
\draw[muted,-latex] (-0.50,0)--(-0.50,0.650);
\draw[teal,thick] (-0.80,0) -- (0.000,0) -- (0.000,0.0000) -- (2.000,0.5000) -- (4.000,0.0000) -- (4.000,0) -- (4.50,0);
\draw[muted] (0.000,0.03)--(0.000,-0.03) node[below,font=\scriptsize,text=navy]{$0$};
\draw[muted] (1.000,0.03)--(1.000,-0.03) node[below,font=\scriptsize,text=navy]{$1$};
\draw[muted] (2.000,0.03)--(2.000,-0.03) node[below,font=\scriptsize,text=navy]{$2$};
\draw[muted] (3.000,0.03)--(3.000,-0.03) node[below,font=\scriptsize,text=navy]{$3$};
\draw[muted] (4.000,0.03)--(4.000,-0.03) node[below,font=\scriptsize,text=navy]{$4$};
\draw[muted,dashed] (-0.50,0.5000)--(4.000,0.5000); \node[left,font=\scriptsize,text=navy] at (-0.50,0.5000) {$\frac12$};
\end{tikzpicture}\end{center}

\end{tasks}

\ex{2} \begin{tasks}(1)
  \task \rest{No}: el área es $3\cdot0{,}5=1{,}5\neq1$.
  \task \rest{No}: toma valores negativos (por ejemplo, $f(0)=-1$).
  \task \rest{Sí}: es positiva y el área es $4\cdot0{,}25=1$ (es la uniforme $U(1,5)$).
\end{tasks}


% =====================================================
\seccio{4}{La distribución normal. La normal tipificada $N(0,1)$ y su tabla}

\exf{1} \begin{tasks}(2)
  \task $\res{0{,}9370}$
  \task $1-0{,}9798=\res{0{,}0202}$
  \task $1-0{,}7939=\res{0{,}2061}$
  \task $P(Z\leq1{,}2)=\res{0{,}8849}$
  \task $0{,}9332-0{,}7734=\res{0{,}1598}$
  \task $0{,}8413-(1-0{,}9772)=\res{0{,}8185}$
  \task $2\cdot0{,}9951-1=\res{0{,}9902}$
\end{tasks}

\ex{1} $2\cdot0{,}8413-1=\res{0{,}6826}$; \ $2\cdot0{,}9772-1=\res{0{,}9544}$; \ $2\cdot0{,}9987-1=\res{0{,}9974}$. \ Entre $\mu\pm\sigma$ hay cerca del 68\,\% de los valores; entre $\mu\pm2\sigma$, el 95\,\%, y entre $\mu\pm3\sigma$, casi todos (99,7\,\%).

\begin{center}\begin{tikzpicture}[x=0.600cm,y=1.6cm,baseline=(current bounding box.center)]
\fill[accent!35] (-1.000,0) -- plot[domain=-1.000:1.000,samples=60,smooth] (\x,{exp(-\x*\x/2)}) -- (1.000,0) -- cycle;
\draw[muted,-latex] (-3.6,0)--(3.7,0);
\draw[teal,thick] plot[domain=-3.5:3.5,samples=90,smooth] (\x,{exp(-\x*\x/2)});
\draw[muted] (-1.000,0.05)--(-1.000,-0.05) node[below,font=\scriptsize,text=navy]{$-1$};
\draw[muted] (0.000,0.05)--(0.000,-0.05) node[below,font=\scriptsize,text=navy]{$0$};
\draw[muted] (1.000,0.05)--(1.000,-0.05) node[below,font=\scriptsize,text=navy]{$1$};
\end{tikzpicture}\qquad\begin{tikzpicture}[x=0.600cm,y=1.6cm,baseline=(current bounding box.center)]
\fill[accent!35] (-2.000,0) -- plot[domain=-2.000:2.000,samples=60,smooth] (\x,{exp(-\x*\x/2)}) -- (2.000,0) -- cycle;
\draw[muted,-latex] (-3.6,0)--(3.7,0);
\draw[teal,thick] plot[domain=-3.5:3.5,samples=90,smooth] (\x,{exp(-\x*\x/2)});
\draw[muted] (-2.000,0.05)--(-2.000,-0.05) node[below,font=\scriptsize,text=navy]{$-2$};
\draw[muted] (0.000,0.05)--(0.000,-0.05) node[below,font=\scriptsize,text=navy]{$0$};
\draw[muted] (2.000,0.05)--(2.000,-0.05) node[below,font=\scriptsize,text=navy]{$2$};
\end{tikzpicture}\end{center}


\ex{2} \begin{tasks}(1)
  \task El valor está en la tabla: $\res{k=1{,}23}$
  \task $P(Z\leq k)=0{,}9495\Rightarrow\res{k=1{,}64}$
  \task $P(Z\leq-k)=0{,}75$: más próximo $0{,}7486\to0{,}67$, así que $\res{k\approx-0{,}67}$
  \task $2P(Z\leq k)-1=0{,}9\Rightarrow P(Z\leq k)=0{,}95$: entre $1{,}64$ y $1{,}65$, $\res{k\approx1{,}645}$
\end{tasks}

\ex{2} $2\,(1-0{,}9332)=\res{0{,}1336}$; \ $0{,}6700-0{,}5=\res{0{,}1700}$; \ $P(Z\leq2{,}35)=\res{0{,}9906}$.

\begin{center}\begin{tikzpicture}[x=0.600cm,y=1.6cm,baseline=(current bounding box.center)]
\fill[accent!35] (-3.500,0) -- plot[domain=-3.500:-1.500,samples=60,smooth] (\x,{exp(-\x*\x/2)}) -- (-1.500,0) -- cycle;
\fill[accent!35] (1.500,0) -- plot[domain=1.500:3.500,samples=60,smooth] (\x,{exp(-\x*\x/2)}) -- (3.500,0) -- cycle;
\draw[muted,-latex] (-3.6,0)--(3.7,0);
\draw[teal,thick] plot[domain=-3.5:3.5,samples=90,smooth] (\x,{exp(-\x*\x/2)});
\draw[muted] (-1.500,0.05)--(-1.500,-0.05) node[below,font=\scriptsize,text=navy]{$-1{,}5$};
\draw[muted] (1.500,0.05)--(1.500,-0.05) node[below,font=\scriptsize,text=navy]{$1{,}5$};
\end{tikzpicture}\qquad\begin{tikzpicture}[x=0.600cm,y=1.6cm,baseline=(current bounding box.center)]
\fill[accent!35] (-0.440,0) -- plot[domain=-0.440:0.000,samples=60,smooth] (\x,{exp(-\x*\x/2)}) -- (0.000,0) -- cycle;
\draw[muted,-latex] (-3.6,0)--(3.7,0);
\draw[teal,thick] plot[domain=-3.5:3.5,samples=90,smooth] (\x,{exp(-\x*\x/2)});
\draw[muted] (-0.440,0.05)--(-0.440,-0.05) node[below,font=\scriptsize,text=navy]{$-0{,}44$};
\draw[muted] (0.000,0.05)--(0.000,-0.05) node[below,font=\scriptsize,text=navy]{$0$};
\end{tikzpicture}\qquad\begin{tikzpicture}[x=0.600cm,y=1.6cm,baseline=(current bounding box.center)]
\fill[accent!35] (-2.350,0) -- plot[domain=-2.350:3.500,samples=60,smooth] (\x,{exp(-\x*\x/2)}) -- (3.500,0) -- cycle;
\draw[muted,-latex] (-3.6,0)--(3.7,0);
\draw[teal,thick] plot[domain=-3.5:3.5,samples=90,smooth] (\x,{exp(-\x*\x/2)});
\draw[muted] (-2.350,0.05)--(-2.350,-0.05) node[below,font=\scriptsize,text=navy]{$-2{,}35$};
\draw[muted] (0.000,0.05)--(0.000,-0.05) node[below,font=\scriptsize,text=navy]{$0$};
\end{tikzpicture}\end{center}



% =====================================================
\seccio{5}{Tipificación. Problemas con la normal}

\exf{1} \begin{tasks}(1)
  \task $P(X\leq25)=P\left(Z\leq\frac{25-15}{10}\right)=P(Z\leq1)=\res{0{,}8413}$
  \task $P(-1\leq Z\leq1)=2\cdot0{,}8413-1=\res{0{,}6826}$
\end{tasks}

\begin{center}\begin{tikzpicture}[x=0.600cm,y=1.6cm,baseline=(current bounding box.center)]
\fill[accent!35] (-1.000,0) -- plot[domain=-1.000:1.000,samples=60,smooth] (\x,{exp(-\x*\x/2)}) -- (1.000,0) -- cycle;
\draw[muted,-latex] (-3.6,0)--(3.7,0);
\draw[teal,thick] plot[domain=-3.5:3.5,samples=90,smooth] (\x,{exp(-\x*\x/2)});
\draw[muted] (-1.000,0.05)--(-1.000,-0.05) node[below,font=\scriptsize,text=navy]{$-1$};
\draw[muted] (0.000,0.05)--(0.000,-0.05) node[below,font=\scriptsize,text=navy]{$0$};
\draw[muted] (1.000,0.05)--(1.000,-0.05) node[below,font=\scriptsize,text=navy]{$1$};
\end{tikzpicture}\end{center}


\ex{1} \begin{tasks}(1)
  \task $P(Z>0{,}67)=1-0{,}7486=\res{0{,}2514}$
  \task $P(Z<-1{,}33)=1-0{,}9082=\res{0{,}0918}$
  \task $P(-0{,}67< Z<3{,}33)=0{,}9996-0{,}2514=\res{0{,}7482}$
\end{tasks}

\ex{1} $P(X>130)=P(Z>2)=1-0{,}9772=0{,}0228$; \ $2000\cdot0{,}0228=45{,}6\approx\res{46}$ alumnos. \ $P(85\leq X\leq115)=P(-1\leq Z\leq1)=0{,}6826$; \ $2000\cdot0{,}6826\approx\res{1365}$ alumnos.

\ex{2} \begin{tasks}(1)
  \task 1,76 es la media de los hombres: $\res{0{,}5}$
  \task $P\left(Z\geq\frac{1{,}76-1{,}62}{0{,}11}\right)=P(Z\geq1{,}27)=1-0{,}8980=\res{0{,}1020}$
  \task Hombre: $0{,}5$; mujer: $1-0{,}102=0{,}898$. \rest{Es más probable que una mujer mida menos de 1,76 m.}
\end{tasks}

\ex{2} Multiplicamos cada probabilidad por el total de empleados:\begin{tasks}(1)
  \task $P(Z>1)=0{,}1587$; \ $250\cdot0{,}1587\approx\res{40}$
  \task $P(Z<-0{,}22)=1-0{,}5871=0{,}4129$; \ $250\cdot0{,}4129\approx\res{103}$
  \task $P(18\leq X<30)=P(-1{,}44\leq Z<-0{,}78)=(1-0{,}7823)-(1-0{,}9251)=0{,}1428$; \ $250\cdot0{,}1428\approx\res{36}$
\end{tasks}

\ex{2} \begin{tasks}(1)
  \task $P(-0{,}38< Z<2)=0{,}9772-(1-0{,}6480)=\res{0{,}6252}$: el $62{,}52\,\%$.
  \task $P(X\geq k)=0{,}1\Rightarrow P(Z\leq z)=0{,}9$: más próximo $0{,}8997\to z=1{,}28$; \ $k=116+1{,}28\cdot42=\res{169{,}76}$ minutos.
\end{tasks}

\begin{center}\begin{tikzpicture}[x=0.600cm,y=1.6cm,baseline=(current bounding box.center)]
\fill[accent!35] (1.280,0) -- plot[domain=1.280:3.500,samples=60,smooth] (\x,{exp(-\x*\x/2)}) -- (3.500,0) -- cycle;
\draw[muted,-latex] (-3.6,0)--(3.7,0);
\draw[teal,thick] plot[domain=-3.5:3.5,samples=90,smooth] (\x,{exp(-\x*\x/2)});
\draw[muted] (0.000,0.05)--(0.000,-0.05) node[below,font=\scriptsize,text=navy]{$0$};
\draw[muted] (1.280,0.05)--(1.280,-0.05) node[below,font=\scriptsize,text=navy]{$1{,}28$};
\end{tikzpicture}\end{center}


\ex{2} Son los cuartiles: $z=\pm0{,}67$ (para $0{,}75$) y $z=0$. \ $163-0{,}67\cdot7=\res{158{,}31}$, \ $\res{163}$ \ y \ $163+0{,}67\cdot7=\res{167{,}69}$ cm.

\ex{2} \begin{tasks}(1)
  \task $P(Z\leq z)=0{,}6$: más próximo $0{,}5987\to z=0{,}25$; \ $250+0{,}25\cdot50=\res{262{,}5}$ g
  \task $P(Z>0{,}6)=1-0{,}7257=0{,}2743$; \ $6000\cdot0{,}2743\approx\res{1646}$ truchas
\end{tasks}

\ex{3} $P\left(Z<\frac{170}{\sigma}\right)=0{,}8023\Rightarrow\frac{170}{\sigma}=0{,}85\Rightarrow\res{\sigma=200}$ euros.

\ex{3} \begin{tasks}(1)
  \task $P\left(Z\leq\frac{183-\mu}{6{,}5}\right)=0{,}88493\approx0{,}8849\Rightarrow\frac{183-\mu}{6{,}5}=1{,}2\Rightarrow\res{\mu=175{,}2}$ cm
  \task Por simetría, $z=-1{,}2$: \ $175{,}2-1{,}2\cdot6{,}5=\res{167{,}4}$ cm
\end{tasks}


% =====================================================
\seccio{6}{Aproximación de la distribución binomial a la normal}

\exf{1} \begin{tasks}(1)
  \task $np=20\geq5$, $nq=80\geq5$: \ $\res{X'\sim N(20,\,4)}$ ($\sigma=\sqrt{100\cdot0{,}2\cdot0{,}8}=4$)
  \task $P(X\leq15)\approx P(X'\leq15{,}5)=P(Z\leq-1{,}13)=1-0{,}8708=\res{0{,}1292}$; \ $P(X=20)\approx P(19{,}5\leq X'\leq20{,}5)=P(-0{,}13\leq Z\leq0{,}13)=\res{0{,}1034}$; \ $P(X\geq25)\approx P(X'\geq24{,}5)=P(Z\geq1{,}13)=\res{0{,}1292}$
\end{tasks}

\ex{1} $X\sim B(400;\ 0{,}5)\approx N(200,\,10)$: \ $P(X<190)=P(X\leq189)\approx P(X'\leq189{,}5)=P(Z\leq-1{,}05)=1-0{,}8531=\res{0{,}1469}$.

\ex{2} $X\sim B(100;\ 0{,}5)\approx N(50,\,5)$: \ $P(X\geq60)\approx P(X'\geq59{,}5)=P(Z\geq1{,}9)=1-0{,}9713=\res{0{,}0287}$.

\begin{center}\begin{tikzpicture}[x=0.600cm,y=1.6cm,baseline=(current bounding box.center)]
\fill[accent!35] (1.900,0) -- plot[domain=1.900:3.500,samples=60,smooth] (\x,{exp(-\x*\x/2)}) -- (3.500,0) -- cycle;
\draw[muted,-latex] (-3.6,0)--(3.7,0);
\draw[teal,thick] plot[domain=-3.5:3.5,samples=90,smooth] (\x,{exp(-\x*\x/2)});
\draw[muted] (0.000,0.05)--(0.000,-0.05) node[below,font=\scriptsize,text=navy]{$0$};
\draw[muted] (1.900,0.05)--(1.900,-0.05) node[below,font=\scriptsize,text=navy]{$1{,}9$};
\end{tikzpicture}\end{center}


\ex{2} \begin{tasks}(1)
  \task $\res{X\sim B(1350;\ 0{,}4)}$
  \task $\mu=540$, $\sigma=\sqrt{324}=18$: \ $P(500\leq X\leq580)\approx P(499{,}5\leq X'\leq580{,}5)=P(-2{,}25\leq Z\leq2{,}25)=2\cdot0{,}9878-1=\res{0{,}9756}$
\end{tasks}

\ex{2} $X\sim B(2500;\ 0{,}1)$: \ $\mu=\res{250}$ accidentes; \ $\sigma=\sqrt{225}=15$. \ $P(X\geq270)\approx P(X'\geq269{,}5)=P(Z\geq1{,}3)=1-0{,}9032=\res{0{,}0968}$.

\ex{3} $np=15$, $nq=35$: \ $N(15;\ \sqrt{10{,}5}\approx3{,}24)$. \ $P(9{,}5\leq X'\leq20{,}5)=P(-1{,}70\leq Z\leq1{,}70)=2\cdot0{,}9554-1=\res{0{,}9108}$. El error es de solo $0{,}0012$.


\end{document}
